\section{Evaluating Pillar 3: Glass-Box Explainability \& HCI}
\label{sec:evaluating-explainability}

This section evaluates the final pillar of MERIT, addressing RQ3. We first verify the mathematical correctness of the Shapley values,
applying the coalition Shapley formalism from explainable ML~\cite{lundberg2017unified} at the data-source level.
Then we conclude with a user study to understand the impact of the explainability layer on recruiter decision-making. In that study, we also
consider usability, confidence calibration and decision reversal rate for features such as Shapley values, audit trail and Blind Mode.
We evaluate mathematical correctness of Shapley attributions and the HCI impact of the Glass-Box layer on recruiters.

\subsection{Study 11: Shapley Value Verification}
\label{subsec:study-11-shapley-verification}

MERIT applies the Shapley formalism at the data-source coalition level 
(Section~\ref{sec:explainable-ai-shapley}), distinct from feature-level SHAP 
explanations for individual model predictions~\cite{lundberg2017unified}.
To ensure source-level attributions are mathematically sound, we verify Efficiency, Null Player, Additivity (per-metric), and Determinism 
on ten synthetic personas ($N = \{\text{CV}, \text{GitHub}, \text{LinkedIn}\}$). We intentionally omit the axiom of Symmetry since source
contributions for MERIT are hetrogenous by design, and therefore cannot produce identical marginal contributions in practice.

\subsubsection{Quantitative Verification Results}
Table~\ref{tab:shapley-verification} reports per-candidate results 
($\sum \phi_i = v(N)$ within floating-point tolerance, negative $\phi_i$ occur when a source triggers a penalty, 
e.g.\ GitHub identity mismatch).

\begin{table}[htbp]
    \centering
    \caption[Study~11: Shapley verification]{Study~11: per-candidate Shapley verification ($\sum \phi_i = v(N)$).}
    \label{tab:shapley-verification}
    \vspace{1.5ex}
    \small
    \setlength{\tabcolsep}{7pt}
    \begin{tabular}{l|ccc|cc}
        \hline
        \textbf{Candidate} & $\phi_{\mathrm{CV}}$ & $\phi_{\mathrm{GH}}$ & $\phi_{\mathrm{LI}}$ & $\sum \phi_i$ & $v(N)$ \\ \hline
        Alex Rivers   & +0.3107 & +0.4607 & +0.0087 & 0.7800 & 0.7800 \\
        Buzz Ward     & +0.3545 &  0.0000 & +0.0075 & 0.3620 & 0.3620 \\
        Carl Corp     & +0.1635 &  0.0000 & +0.0155 & 0.1790 & 0.1790 \\
        Felix Vance   & +0.3735 & +0.0660 & +0.0075 & 0.4470 & 0.4470 \\
        Fiona Frost   & +0.2660 & +0.0660 & +0.0080 & 0.3400 & 0.3400 \\
        Ghost Gary    & +0.2940 &  0.0000 &  0.0000 & 0.2940 & 0.2940 \\
        Jordan Smith  & +0.1028 & +0.5288 & +0.0113 & 0.6430 & 0.6430 \\
        Kim Junior    & +0.1340 & +0.2550 &  0.0000 & 0.3890 & 0.3890 \\
        Sam Old       & +0.1648 & +0.1878 & +0.0153 & 0.3680 & 0.3680 \\
        Vince Vault   & +0.2652 & $-$0.0708 & +0.0117 & 0.2060 & 0.2060 \\ \hline
    \end{tabular}
\end{table}

\begin{itemize}
    \item \textbf{Efficiency:} For all 10 candidates, $\sum \phi_i = v(N)$ within floating-point tolerance (Table~\ref{tab:shapley-verification}).

    \item \textbf{Null player:} Ghost Gary receives $\phi_{\text{GitHub}} = \phi_{\text{LinkedIn}} = 0$ on empty channels.

    \item \textbf{Determinism:} Five repeated runs produced zero drift (exact analytical Shapley, not Monte Carlo).

    \item \textbf{Per-metric efficiency:} Per-metric Shapley sums matched grand-coalition metric scores across all candidate--metric pairs.
\end{itemize}

For a visualisation of these results, refer to Figure~\ref{fig:shapley-combined} in 
Appendix~\ref{appendix:shapley-charts}.

\subsubsection{Key XAI Insights}

One key finding from this study is visible in the GitHub Shapley value for Vince Vault, where we show the presence of negative attributions.
In the context of cooperative game theory, this would occur when the introduction of a player leads to a reduction in the overall coalition's
winnings. For MERIT, this would occur when the introduction of a source leads to a fall in that candidate's match score. In this case,
it was due to the introduction of malicious intent via identity squatting, leading to a significant penalty being applied to the candidate's score.
Since Felix only provided false evidence for the GitHub source, negative attributions were focused there, rather than on LinkedIn or CV.
This penalty is then able to be presented to the recruiter to show that a candidate is allegedly claiming evidence that they do not have.
It is now the recruiter's responsibility to investigate the engine further to see if the fuzzy match failed due to actual intent of
malice, or due to a false positive.

Another finding can be presented through our hidden gem in Section~\ref{subsec:study-01-comparative-ablation},
Jordan Smith. As shown in that study, pivot from MERIT (CV Only) to MERIT (Full) boosts his score from rank 10 to rank 2. Due
to the introduction of Shapley values, we can truly confirm that this boost in rank is due to the inclusion of the GitHub data source,
which contributed $\phi_{GitHub} = 0.529$ to his score. This is presented to the recruiter on the frontend to show that the majority
of his score is demonstrated, not claimed. Once again, it is now the recruiter's decision to investigate the applicant further to
see if they are a good fit for the role, since they have demonstrated well that they align with the job, but their claims don't
seem to carry the same weight as their demonstrated evidence.

It is important to note that the dataset for this study reuses that of Study~01, therefore, we do not have significant Shapley values for
any candidate when it comes to the LinkedIn source attribution. This is due to the study being designed to evaluate the effect of
demonstrated evidence when it comes to candidate selection, so we kept LinkedIn evidence minimal to stress-test this CV-versus-GitHub contrast.
See Section~\ref{subsec:construct-validity} for more details. Note that this does not affect
the validity of this study in any way, we are primarily providing an explanation as to why LinkedIn doesn't contain significant
Shapley values.

\begin{table}[htbp]
    \centering
    \caption[Study~11: Aggregate Shapley by source]{Study~11: aggregate positive Shapley shares by source.}
    \label{tab:shapley-source-dominance}
    \vspace{1.5ex}
    \small
    \begin{tabular}{l|c|c}
        \hline
        \textbf{Source} & \textbf{Sum of positive $\phi_i$} & \textbf{Share of pool (\%)} \\ \hline
        CV (claimed)       & 2.429 & 59.6 \\
        GitHub (demonstrated) & 1.564 & 38.4 \\
        LinkedIn (claimed)    & 0.085 & 2.1 \\ \hline
    \end{tabular}
\end{table}

We can observe that in this cohort, CV evidence dominates the overall score, contributing to 59.6\% of the overall score of 
the candidates. This is followed by demonstrated GitHub evidence, contributing to 38.4\% of the overall score.
While LinkedIn evidence is present, it contributes to only 2.1\% of the overall score, indicating that it is not a significant
factor in the candidate's score. Note the margin for error since we only consider positive $\phi_i$ for the denominator (we exclude 
the negative GitHub attribution from Felix Vance).

\subsection{Study 14: HCI Trial}
\label{subsec:hci-trial-protocol}

Study 11 mathematically confirms that the explainable AI (XAI) engine is consistent, Study 14 evaluates the human-facing 
glass-box design against the transparency deficit identified in algorithmic hiring reviews~\cite{bogen2018,raghavan2020}.
Lo et al.~\cite{lo2025aihiring} report HR agreement on generated scores, but their published evaluation does not include a 
before/after test of whether recruiters reverse shortlist decisions once explanations are shown. 
We include the Decision Reversal Rate (DRR) metric to target this gap.
The purpose of our XAI is to provide transparent explanations to the recruiter, so we deem it important to evaluate
how our Glass-Box influences recruiter behaviour, hence our need for this metric.

In this study, we measure not only the decision, but also the confidence of recruiters with regards to their decision.
We also measure how usable this engine is, since if the system is too difficult to use, recruiters would likely avoid using it,
falling back into automation bias to avoid significant cognitive load. Therefore, we have the following hypotheses:

\begin{itemize}
    \item \textbf{Hypothesis 1 (Decision Calibration):} Access to the Glass-Box XAI audit trail will increase the rate at 
    which recruiters reverse incorrect screening decisions for adversarial candidates (such as
          keyword-stuffers and identity squatters) and "hidden gems" (like high-potential candidates with low tenure).
    \item \textbf{Hypothesis 2 (Confidence Calibration):} Recruiters will exhibit a higher alignment between their
          screening confidence and the empirical evidence, resulting in increased confidence for verified candidates and
          decreased confidence for unverified profiles.
    \item \textbf{Hypothesis 3 (Usability and Trust):} The interactive explainability layer will achieve a high System
          Usability Scale (SUS) score and receive high qualitative ratings for transparency and defensibility compared to a
          traditional black-box score.
\end{itemize}

\subsubsection{Two-Stage HCI Trial}
We start with a two-stage user study involving professional technical recruiters and hiring managers as participants.
This trial follows a two-stage process for each candidate in the pool:

\subsubsection{Stage 1: Control Condition (Black-Box Score)}
Recruiters are given a candidate CV with a single aggregate match alongside per-metric scores, but with no explanations.
Recruiters must make a decision on whether to shortlist or reject the candidate, and rate their confidence on a Likert scale
(1 = Not Confident at all, 5 = Extremely Confident). They must also provide a brief qualitative justification for their decision.

\subsubsection{Stage 2: Treatment Condition (Glass-Box Scorecard)}
Recruiters are now shown the same candidate, but with access to the glass-box scorecards, including the Shapley layer, audit trail,
and other explainability features. Recruiters will then be asked to state whether they wish to reverse or maintain their decision.
They are also asked to rate their confidence, along with an explanation of how the visualisations influenced their assessment, if at all.

\subsubsection{Key Quantitative Metrics and Empirical Results}
This pilot user trial includes 5 participants, consisting of 2 professional recruiters ($R_1, R_2$) with extensive experience in
software development talent acquisition, and 3 technical peers ($P_1, P_2, P_3$), MEng Computing students at Imperial College London, all of whom were acquainted with the author. 
These participants are familiar with the technical recruitment process and have undergone multiple technical projects and internships. 
Given $N=5$, we treat this as a formative pilot and report descriptive statistics, inferential tests are 
exploratory only (Section~\ref{sec:threats-to-validity}).

To prevent participant fatigue, we present each evaluator with only 4 synthetic candidates to evaluate. Unlike the personas
mentioned in previous sections, these candidates were designed in a real-world scenario. We used real candidate information
as the basis for synthesising these candidates:
\begin{itemize}
    \item \textbf{Moinecha Ramos-Horta} (Verifiable Elite Developer)
    \item \textbf{Yusuf Nilsson} (The "Hidden Gem" - 1 year of experience on CV, elite GitHub)
    \item \textbf{Robert Todd} (Adversarial Keyword Stuffer)
    \item \textbf{Aaron Mohammed} (Identity Squatter - claims deep Java/Go systems expertise on CV, squatted GitHub)
\end{itemize}

We present the quantitative results of this trial in Table~\ref{tab:hci-trial-results}.
Stage~1 and Stage~2 counts use $S$ (shortlist) and $R$ (reject), \textbf{DRR} is the fraction of Stage~1 decisions reversed after the Glass-Box scorecard.
Confidence uses a 5-point Likert scale (1 = not confident, 5 = extremely confident).
Decision-reversal, confidence-shift, and SUS distributions are in Figure~\ref{fig:hci-combined}
(panels a--c, Appendix~\ref{appendix:hci-trial-charts}). 
Raw session data are in \texttt{evaluation/14-hci\_trial/}.


\begin{table}[htbp]
    \centering
    \caption[Study~14: HCI pilot results]{Study~14: HCI pilot results ($N = 5$ evaluators).}
    \label{tab:hci-trial-results}
    \vspace{1.5ex}
    \small
    \setlength{\tabcolsep}{5pt}
    \begin{tabular}{l|c|c|c|c|c|c}
        \hline
        \textbf{Candidate}   & \textbf{Stage 1 (ctrl)} & \textbf{Stage 2 (trt)} & \textbf{DRR} & \textbf{$C_1$ (1--5)} & \textbf{$C_2$ (1--5)} & \textbf{$\Delta C$} \\ \hline
        Moinecha Ramos-Horta & $5S$ / $0R$            & $5S$ / $0R$            & 0.0\%        & 3.80                    & 4.60                    & +0.80                 \\
        Yusuf Nilsson        & $0S$ / $5R$            & $4S$ / $1R$            & 80.0\%       & 4.20                    & 4.80                    & +0.60                 \\
        Robert Todd          & $4S$ / $1R$            & $1S$ / $4R$            & 75.0\%       & 3.60                    & 4.40                    & +0.80                 \\
        Aaron Mohammed       & $3S$ / $2R$            & $0S$ / $5R$            & 100.0\%      & 3.20                    & 4.80                    & +1.60                 \\ \hline
    \end{tabular}
\end{table}

This pilot provides directional support for all three hypotheses under realistic hiring constraints, 
none are confirmed at inferential strength given $N=5$. We detail the results below:

\subsection*{Quantitative Results}

\subsubsection{Decision Calibration (Hypothesis 1)}
During the first stage, we can see the issue of the ``resume black hole'': all 5 evaluators rejected Yusuf Nilsson
due to his low tenure (1 year of experience on CV). However, upon transitioning to stage 2, we see that 4/5 evaluators
(80\% DRR) reversed their decision, recognising his valuable GitHub contributions. Recruiter $R_2$ chose to
maintain the rejection, noting that the JD strictly required 3+ years of experience regardless
of technical strength. This illustrates how human decision-makers calibrate system outputs with non-quantifiable
operational constraints.

Another key finding is that the explainability layer successfully aided in unmasking adversarial manipulation. More specifically,
we turn to the candidate \textbf{Robert Todd}, who managed to fool 4 out of 5 evaluators in stage 1 due to his keyword-dense
resume and aggregate score. However, upon transitioning to stage 2, we see that 3 out of 4 of those evaluators (75\% DRR)
reversed their decision, recognising his keyword stuffing.

The most significant instance of this is with the candidate \textbf{Aaron Mohammed}, who managed to fool 3 out of 5 evaluators
in stage 1 due to his gamified CV. However, upon transitioning to stage 2, rather than seeing supported evidence, recruiters were
presented with a significant red flag. Aaron Mohammed provided a GitHub profile that was not their own in order to pass off
someone else's work as his own. This is a serious ethical breach and a violation of professional integrity. Due to this, all
3 evaluators reversed their decision, rejecting Aaron Mohammed. It is crucial to note that recruiters did not reject
Mohammed based on the fact that he was not a good developer (in fact, according to his CV he was an expert in Java and Go).
Instead, they rejected him based on the fact that he showed a lack of integrity and dishonesty. This illustrates how the explainability layer
can be used to identify and reject participants who are not a good fit for the company culture, even if they are technically skilled.


\subsubsection{Confidence Calibration (Hypothesis 2)}
Once evaluators were presented with the Glass-Box MERIT interface, we observed a positive confidence shift ($\Delta C$)
across all profiles. This was particularly notable in the case of Moinecha Ramos-Horta, where her average confidence
rose from $3.80$ (Stage 1) to $4.60$ (Stage 2), an increase of $0.80$. In a similar fashion, confidence for Aaron Mohammed
also increased, rising from $3.20$ (Stage 1) to $4.80$ (Stage 2), an increase of $1.60$, which was
due to the veto penalty that stemmed from his dishonesty, giving the evaluators the confidence to reject him with high conviction.

The average confidence across all 20 profile-level observations (5 participants $\times$ 4 candidates) increased from $3.75$ 
(Stage 1) to $4.45$ (Stage 2), a mean increase of $0.70$ on a 5-point scale. A paired $t$-test on these rows yields 
$t(19) = -5.60$, $p < 0.001$, but the rows are not independent (fixed Stage~1$\rightarrow$Stage~2 order, repeated participants, 
acquaintance bias among peers). We therefore report this as descriptive evidence for Hypothesis~2, not confirmatory inference.

\subsubsection{Usability (Hypothesis 3)}
To evaluate the usability of MERIT's glass-box architecture, we followed the standard System Usability Scale (SUS) survey \cite{brooke1996sus}.
MERIT achieved an average SUS score of \textbf{81.50} out of 100, with individual scores ranging from 77.5 to 85.0.
According to the Bangor et al. adjective rating scale for SUS scores \cite{bangor2009_sus}, this corresponds to a ``Good'' rating (Grade B) at face value.
With $N=5$, it indicates interface feasibility only, not product-market validation (Section~\ref{subsec:construct-validity}).

\subsubsection{Qualitative themes}

The justifications presented by our evaluators can be found in \texttt{evaluation/14-hci\_trial/test\_data/responses.json}. As shown, there are
alignments with the quantitative results in Table~\ref{tab:hci-trial-results}. Take the evaluator
that still shortlisted Robert Todd after the keyword-stuffing alert. Despite what is considered a ``serious red flag'' by $P_3$, they still plan to 
``keep him shortlisted'' for a brief phone screen. Note how the implementation of the glass-box not only introduces contestability with the 
scores presented, but also offers recruiters the means to investigate claims made by candidates in better detail where evidence may be missing.

Meanwhile, other penalties such as identity squatting are treated as decisive when malicious dishonesty is surfaced. In the case of Aaron Mohammed,
recruiter $R_1$ states that this is an ``immediate reject'' due to the ``Identity Mismatch Veto'' penalty upon seeing that the GitHub account's
full name does not match the candidate's name. Passing off another's work as one's own is a serious breach of honesty and integrity, and is often 
not a good fit for company culture. 

The responses point towards cases in which automation bias can be challenged by a recruiter's human judgement. While this study confirms the 
theoretical success of the Glass-Box design, we must note that this is a pilot study. Studies with a larger recruiter panel would be needed
before stronger claims can be made.

\textbf{Outcome:} Meets the internal SUS target (81.5/100).
Decision-reversal and confidence shifts provide directional support for glass-box calibration, but the pilot is not confirmatory ($N=5$).